\documentclass[11pt]{article}
\usepackage[utf8]{inputenc}
\usepackage[T1]{fontenc}
\usepackage{amsmath,amssymb,mathtools}
\usepackage[margin=1in]{geometry}
\usepackage{hyperref}
\usepackage{array}
\usepackage{parskip}
\setlength{\parindent}{0pt}
\hypersetup{colorlinks=true,linkcolor=blue,urlcolor=blue}
\title{Lecture 12: Fourier Transforms (CT and DT)}
\author{}
\date{}
\begin{document}
\maketitle

\section*{From Fourier Series to Fourier Transform}

An aperiodic signal is a periodic signal with $T \to \infty$:
\begin{itemize}
  \item Harmonics $k\omega_0$ become infinitesimally close ($\omega_0 \to 0$)
  \item Discrete spectrum becomes a \textbf{continuous spectrum}
  \item Sum becomes an integral
\end{itemize}

\section*{CT Fourier Transform (CTFT) Pair}

\subsection*{Forward Transform (Analysis): time -> frequency}

\[X(j\omega) = \int_{-\infty}^{\infty} x(t) \, e^{-j\omega t} \, dt\]

\subsection*{Inverse Transform (Synthesis): frequency -> time}

\[x(t) = \frac{1}{2\pi} \int_{-\infty}^{\infty} X(j\omega) \, e^{j\omega t} \, d\omega\]

\subsection*{Convergence}

Same Dirichlet conditions as CTFS, plus absolute integrability:

\[\int_{-\infty}^{\infty} |x(t)| \, dt < \infty\]

\section*{Basic CT FT Pairs}

\begin{center}
\begin{tabular}{|l|l|}
\hline
\textbf{Signal $x(t)$} & \textbf{Transform $X(j\omega)$} \\
\hline
$e^{-at}u(t), \; a > 0$ & $\dfrac{1}{a + j\omega}$ \\
\hline
$e^{-a\|t\|}, \; a > 0$ & $\dfrac{2a}{a^2 + \omega^2}$ \\
\hline
$\delta(t)$ & $1$ \\
\hline
$1$ & $2\pi\delta(\omega)$ \\
\hline
$e^{j\omega_0 t}$ & $2\pi\delta(\omega - \omega_0)$ \\
\hline
$\cos(\omega_0 t)$ & $\pi[\delta(\omega - \omega_0) + \delta(\omega + \omega_0)]$ \\
\hline
$\sin(\omega_0 t)$ & $\frac{\pi}{j}[\delta(\omega - \omega_0) - \delta(\omega + \omega_0)]$ \\
\hline
Rectangular pulse (width $T_1$) & $\dfrac{2\sin(\omega T_1)}{\omega} = 2T_1 \text{sinc}\left(\frac{\omega T_1}{\pi}\right)$ \\
\hline
\end{tabular}
\end{center}

\section*{FT of Periodic Signals}

A periodic signal $x(t) = \sum a_k e^{jk\omega_0 t}$ has FT:

\[X(j\omega) = \sum_{k=-\infty}^{\infty} 2\pi a_k \, \delta(\omega - k\omega_0)\]

The spectrum is a \textbf{train of impulses} at harmonic frequencies, weighted by $2\pi a_k$.

\section*{Discrete-Time Fourier Transform (DTFT)}

\subsection*{DTFT Pair}

\textbf{Forward:}
\[X(e^{j\omega}) = \sum_{n=-\infty}^{\infty} x[n] \, e^{-j\omega n}\]

\textbf{Inverse:}
\[x[n] = \frac{1}{2\pi} \int_{2\pi} X(e^{j\omega}) \, e^{j\omega n} \, d\omega\]

\subsection*{Key DTFT Properties}
\begin{itemize}
  \item $X(e^{j\omega})$ is always \textbf{periodic in $\omega$} with period $2\pi$
  \item Only need to specify over $[0, 2\pi)$ or $[-\pi, \pi)$
\end{itemize}

\subsection*{Basic DT FT Pairs}

\begin{center}
\begin{tabular}{|l|l|}
\hline
\textbf{Signal $x[n]$} & \textbf{Transform $X(e^{j\omega})$} \\
\hline
$a^n u[n], \; \|a\| < 1$ & $\dfrac{1}{1 - ae^{-j\omega}}$ \\
\hline
$\delta[n]$ & $1$ \\
\hline
$1$ & $2\pi \sum_k \delta(\omega - 2\pi k)$ \\
\hline
$e^{j\omega_0 n}$ & $2\pi \sum_k \delta(\omega - \omega_0 - 2\pi k)$ \\
\hline
\end{tabular}
\end{center}

\section*{The Four Fourier Representations}

\begin{center}
\begin{tabular}{|l|l|l|l|}
\hline
\textbf{Signal Type} & \textbf{Time} & \textbf{Frequency} & \textbf{Representation} \\
\hline
CT Periodic & Continuous & Discrete & CTFS \\
\hline
CT Aperiodic & Continuous & Continuous & CTFT \\
\hline
DT Periodic & Discrete & Discrete & DTFS \\
\hline
DT Aperiodic & Discrete & Continuous & DTFT \\
\hline
\end{tabular}
\end{center}

\textbf{Pattern:} Discrete in one domain $\Leftrightarrow$ Periodic in the other domain.

\end{document}
